% NLA Architecture: AV / Actor — Activation → Text
% % NLA Architecture: AV / Actor — Activation → Text
% % NLA Architecture: AV / Actor — Activation → Text
% % NLA Architecture: AV / Actor — Activation → Text
% \input{latex/02-av-actor}

% ============================================================
% Section 1: Architecture
% ============================================================

\begin{frame}{AV / Architecture — Core Design}
\small

\textbf{Activation Verbalizer (AV)} 将残差流 activation 翻译为自然语言 explanation。AV 和 AR 共同构成 autoencoder：AV 编码 activation $\rightarrow$ text，AR 解码 text $\rightarrow$ activation。

\begin{block}{架构}
\begin{itemize}
  \item 初始化为目标 LLM 副本（Qwen2.5-7B-Instruct）;构造 prompt：\code{"Explain: <concept><INJECT></concept>"}
  \item Activation 通过特殊 token \code{※} (U+203B) 注入。将该位置的 \textbf{embedding 向量} 替换为保存的 activation vector（3584-d）。
  \item LLM 从注入位置继续 decode，生成自然语言 explanation。
  \item 输出包裹在 \code{<explanation>...</explanation>} 标签中。
\end{itemize}
\end{block}

\end{frame}


\begin{frame}{AV / Architecture — Injection Mechanism \& I/O Format}
\small

\begin{itemize}
  \item Hook 只修改前向传播中的数据流（替换一个 token 的 embedding），不改变任何参数。
  \item 不在目标层（layer 20）注入：避免与 gradient checkpointing / FSDP 的交互复杂性。
  \item Embedding 层是唯一在所有模型架构中都稳定可 hook 的点。
  \item Activation 虽提取自 layer 20，但从 layer 0 开始 decode 给模型更多「理解」这个向量的空间。
\end{itemize}

\vspace{0.4em}
\begin{center}
\begin{tikzpicture}[
  node distance=1.0cm,
  box/.style={draw, rounded corners=2pt, align=center, minimum width=1.6cm, minimum height=0.65cm, font=\footnotesize},
  arr/.style={->, thick, color=nlablue}
]
\node[box, fill=blue!8] (tok) {Token 序列\\{\footnotesize \code{[B, S]}}};
\node[box, fill=green!8, right=0.5cm of tok] (emb) {Embedding\\{\footnotesize \code{[B,S,3584]}}};
\node[box, fill=orange!8, right=0.5cm of emb] (inj) {Hook 替换\\{\footnotesize ※ 位置}};
\node[box, fill=blue!8, right=0.5cm of inj] (tr) {28 层\\Transformer};
\node[box, fill=blue!8, right=0.5cm of tr] (head) {LM Head};
\node[box, fill=red!8, right=0.5cm of head] (out) {Logits\\{\footnotesize \code{[B,S,V]}}};
\node[box, fill=green!8, above=0.7cm of inj] (act) {Activation\\{\footnotesize \code{[B,3584]}}};

\draw[arr] (tok) -- (emb);
\draw[arr] (emb) -- (inj);
\draw[arr] (inj) -- (tr);
\draw[arr] (tr) -- (head);
\draw[arr] (head) -- (out);
\draw[arr] (act) -- (inj);
\end{tikzpicture}
\end{center}
\end{frame}


% ============================================================
% Section 2: Training Method
% ============================================================

\begin{frame}{AV / Training：SFT + RL}
\small

AV 训练分两阶段：先用 API 生成的 explanation 做 warm-start，再用 RL 以重建保真度为 reward 继续优化。

\begin{columns}[T,onlytextwidth]
\begin{column}{0.48\linewidth}
\begin{block}{Stage 1a: AV-SFT（监督微调）}
\begin{itemize}
  \item 用 GPT-4 / Claude 等 API 为 activation 生成 reference explanation。
  \item 训练 AV 模彷 API 的解释风格和格式。
  \item 数据格式：prompt 含 \code{<INJECT>} token，response 为 \code{<explanation>...\\</explanation>}。
  \item Warm-start 让 AV 学会基本的 explanation 格式，为后续 RL 提供合理的初始策略。
\end{itemize}
\end{block}
\end{column}
\begin{column}{0.48\linewidth}
\begin{block}{Stage 1b: AV-RL（强化学习）}
\begin{itemize}
  \item 用 \textbf{GRPO} 继续训练。
  \item Reward = \textbf{AR cosine similarity} + \textbf{brevity bonus}。
  \item 鼓励 AR 能精确重建的紧凑 explanation。
\end{itemize}
\end{block}
\end{column}
\end{columns}

\vspace{0.5em}

Reward 信号来自\textbf{重建保真度}（AR cosine similarity）。但实践中，explanation 随 RL 训练自然地变得更 informative——因为更精确的描述有助于 AR 重建。

\end{frame}

% ============================================================
% Section 1: Architecture
% ============================================================

\begin{frame}{AV / Architecture — Core Design}
\small

\textbf{Activation Verbalizer (AV)} 将残差流 activation 翻译为自然语言 explanation。AV 和 AR 共同构成 autoencoder：AV 编码 activation $\rightarrow$ text，AR 解码 text $\rightarrow$ activation。

\begin{block}{架构}
\begin{itemize}
  \item 初始化为目标 LLM 副本（Qwen2.5-7B-Instruct）;构造 prompt：\code{"Explain: <concept><INJECT></concept>"}
  \item Activation 通过特殊 token \code{※} (U+203B) 注入。将该位置的 \textbf{embedding 向量} 替换为保存的 activation vector（3584-d）。
  \item LLM 从注入位置继续 decode，生成自然语言 explanation。
  \item 输出包裹在 \code{<explanation>...</explanation>} 标签中。
\end{itemize}
\end{block}

\end{frame}


\begin{frame}{AV / Architecture — Injection Mechanism \& I/O Format}
\small

\begin{itemize}
  \item Hook 只修改前向传播中的数据流（替换一个 token 的 embedding），不改变任何参数。
  \item 不在目标层（layer 20）注入：避免与 gradient checkpointing / FSDP 的交互复杂性。
  \item Embedding 层是唯一在所有模型架构中都稳定可 hook 的点。
  \item Activation 虽提取自 layer 20，但从 layer 0 开始 decode 给模型更多「理解」这个向量的空间。
\end{itemize}

\vspace{0.4em}
\begin{center}
\begin{tikzpicture}[
  node distance=1.0cm,
  box/.style={draw, rounded corners=2pt, align=center, minimum width=1.6cm, minimum height=0.65cm, font=\footnotesize},
  arr/.style={->, thick, color=nlablue}
]
\node[box, fill=blue!8] (tok) {Token 序列\\{\footnotesize \code{[B, S]}}};
\node[box, fill=green!8, right=0.5cm of tok] (emb) {Embedding\\{\footnotesize \code{[B,S,3584]}}};
\node[box, fill=orange!8, right=0.5cm of emb] (inj) {Hook 替换\\{\footnotesize ※ 位置}};
\node[box, fill=blue!8, right=0.5cm of inj] (tr) {28 层\\Transformer};
\node[box, fill=blue!8, right=0.5cm of tr] (head) {LM Head};
\node[box, fill=red!8, right=0.5cm of head] (out) {Logits\\{\footnotesize \code{[B,S,V]}}};
\node[box, fill=green!8, above=0.7cm of inj] (act) {Activation\\{\footnotesize \code{[B,3584]}}};

\draw[arr] (tok) -- (emb);
\draw[arr] (emb) -- (inj);
\draw[arr] (inj) -- (tr);
\draw[arr] (tr) -- (head);
\draw[arr] (head) -- (out);
\draw[arr] (act) -- (inj);
\end{tikzpicture}
\end{center}
\end{frame}


% ============================================================
% Section 2: Training Method
% ============================================================

\begin{frame}{AV / Training：SFT + RL}
\small

AV 训练分两阶段：先用 API 生成的 explanation 做 warm-start，再用 RL 以重建保真度为 reward 继续优化。

\begin{columns}[T,onlytextwidth]
\begin{column}{0.48\linewidth}
\begin{block}{Stage 1a: AV-SFT（监督微调）}
\begin{itemize}
  \item 用 GPT-4 / Claude 等 API 为 activation 生成 reference explanation。
  \item 训练 AV 模彷 API 的解释风格和格式。
  \item 数据格式：prompt 含 \code{<INJECT>} token，response 为 \code{<explanation>...\\</explanation>}。
  \item Warm-start 让 AV 学会基本的 explanation 格式，为后续 RL 提供合理的初始策略。
\end{itemize}
\end{block}
\end{column}
\begin{column}{0.48\linewidth}
\begin{block}{Stage 1b: AV-RL（强化学习）}
\begin{itemize}
  \item 用 \textbf{GRPO} 继续训练。
  \item Reward = \textbf{AR cosine similarity} + \textbf{brevity bonus}。
  \item 鼓励 AR 能精确重建的紧凑 explanation。
\end{itemize}
\end{block}
\end{column}
\end{columns}

\vspace{0.5em}

Reward 信号来自\textbf{重建保真度}（AR cosine similarity）。但实践中，explanation 随 RL 训练自然地变得更 informative——因为更精确的描述有助于 AR 重建。

\end{frame}

% ============================================================
% Section 1: Architecture
% ============================================================

\begin{frame}{AV / Architecture — Core Design}
\small

\textbf{Activation Verbalizer (AV)} 将残差流 activation 翻译为自然语言 explanation。AV 和 AR 共同构成 autoencoder：AV 编码 activation $\rightarrow$ text，AR 解码 text $\rightarrow$ activation。

\begin{block}{架构}
\begin{itemize}
  \item 初始化为目标 LLM 副本（Qwen2.5-7B-Instruct）;构造 prompt：\code{"Explain: <concept><INJECT></concept>"}
  \item Activation 通过特殊 token \code{※} (U+203B) 注入。将该位置的 \textbf{embedding 向量} 替换为保存的 activation vector（3584-d）。
  \item LLM 从注入位置继续 decode，生成自然语言 explanation。
  \item 输出包裹在 \code{<explanation>...</explanation>} 标签中。
\end{itemize}
\end{block}

\end{frame}


\begin{frame}{AV / Architecture — Injection Mechanism \& I/O Format}
\small

\begin{itemize}
  \item Hook 只修改前向传播中的数据流（替换一个 token 的 embedding），不改变任何参数。
  \item 不在目标层（layer 20）注入：避免与 gradient checkpointing / FSDP 的交互复杂性。
  \item Embedding 层是唯一在所有模型架构中都稳定可 hook 的点。
  \item Activation 虽提取自 layer 20，但从 layer 0 开始 decode 给模型更多「理解」这个向量的空间。
\end{itemize}

\vspace{0.4em}
\begin{center}
\begin{tikzpicture}[
  node distance=1.0cm,
  box/.style={draw, rounded corners=2pt, align=center, minimum width=1.6cm, minimum height=0.65cm, font=\footnotesize},
  arr/.style={->, thick, color=nlablue}
]
\node[box, fill=blue!8] (tok) {Token 序列\\{\footnotesize \code{[B, S]}}};
\node[box, fill=green!8, right=0.5cm of tok] (emb) {Embedding\\{\footnotesize \code{[B,S,3584]}}};
\node[box, fill=orange!8, right=0.5cm of emb] (inj) {Hook 替换\\{\footnotesize ※ 位置}};
\node[box, fill=blue!8, right=0.5cm of inj] (tr) {28 层\\Transformer};
\node[box, fill=blue!8, right=0.5cm of tr] (head) {LM Head};
\node[box, fill=red!8, right=0.5cm of head] (out) {Logits\\{\footnotesize \code{[B,S,V]}}};
\node[box, fill=green!8, above=0.7cm of inj] (act) {Activation\\{\footnotesize \code{[B,3584]}}};

\draw[arr] (tok) -- (emb);
\draw[arr] (emb) -- (inj);
\draw[arr] (inj) -- (tr);
\draw[arr] (tr) -- (head);
\draw[arr] (head) -- (out);
\draw[arr] (act) -- (inj);
\end{tikzpicture}
\end{center}
\end{frame}


% ============================================================
% Section 2: Training Method
% ============================================================

\begin{frame}{AV / Training：SFT + RL}
\small

AV 训练分两阶段：先用 API 生成的 explanation 做 warm-start，再用 RL 以重建保真度为 reward 继续优化。

\begin{columns}[T,onlytextwidth]
\begin{column}{0.48\linewidth}
\begin{block}{Stage 1a: AV-SFT（监督微调）}
\begin{itemize}
  \item 用 GPT-4 / Claude 等 API 为 activation 生成 reference explanation。
  \item 训练 AV 模彷 API 的解释风格和格式。
  \item 数据格式：prompt 含 \code{<INJECT>} token，response 为 \code{<explanation>...\\</explanation>}。
  \item Warm-start 让 AV 学会基本的 explanation 格式，为后续 RL 提供合理的初始策略。
\end{itemize}
\end{block}
\end{column}
\begin{column}{0.48\linewidth}
\begin{block}{Stage 1b: AV-RL（强化学习）}
\begin{itemize}
  \item 用 \textbf{GRPO} 继续训练。
  \item Reward = \textbf{AR cosine similarity} + \textbf{brevity bonus}。
  \item 鼓励 AR 能精确重建的紧凑 explanation。
\end{itemize}
\end{block}
\end{column}
\end{columns}

\vspace{0.5em}

Reward 信号来自\textbf{重建保真度}（AR cosine similarity）。但实践中，explanation 随 RL 训练自然地变得更 informative——因为更精确的描述有助于 AR 重建。

\end{frame}

% ============================================================
% Section 1: Architecture
% ============================================================

\begin{frame}{AV / Architecture — Core Design}
\small

\textbf{Activation Verbalizer (AV)} 将残差流 activation 翻译为自然语言 explanation。AV 和 AR 共同构成 autoencoder：AV 编码 activation $\rightarrow$ text，AR 解码 text $\rightarrow$ activation。

\begin{block}{架构}
\begin{itemize}
  \item 初始化为目标 LLM 副本（Qwen2.5-7B-Instruct）;构造 prompt：\code{"Explain: <concept><INJECT></concept>"}
  \item Activation 通过特殊 token \code{※} (U+203B) 注入。将该位置的 \textbf{embedding 向量} 替换为保存的 activation vector（3584-d）。
  \item LLM 从注入位置继续 decode，生成自然语言 explanation。
  \item 输出包裹在 \code{<explanation>...</explanation>} 标签中。
\end{itemize}
\end{block}

\end{frame}


\begin{frame}{AV / Architecture — Injection Mechanism \& I/O Format}
\small

\begin{itemize}
  \item Hook 只修改前向传播中的数据流（替换一个 token 的 embedding），不改变任何参数。
  \item 不在目标层（layer 20）注入：避免与 gradient checkpointing / FSDP 的交互复杂性。
  \item Embedding 层是唯一在所有模型架构中都稳定可 hook 的点。
  \item Activation 虽提取自 layer 20，但从 layer 0 开始 decode 给模型更多「理解」这个向量的空间。
\end{itemize}

\vspace{0.4em}
\begin{center}
\begin{tikzpicture}[
  node distance=1.0cm,
  box/.style={draw, rounded corners=2pt, align=center, minimum width=1.6cm, minimum height=0.65cm, font=\footnotesize},
  arr/.style={->, thick, color=nlablue}
]
\node[box, fill=blue!8] (tok) {Token 序列\\{\footnotesize \code{[B, S]}}};
\node[box, fill=green!8, right=0.5cm of tok] (emb) {Embedding\\{\footnotesize \code{[B,S,3584]}}};
\node[box, fill=orange!8, right=0.5cm of emb] (inj) {Hook 替换\\{\footnotesize ※ 位置}};
\node[box, fill=blue!8, right=0.5cm of inj] (tr) {28 层\\Transformer};
\node[box, fill=blue!8, right=0.5cm of tr] (head) {LM Head};
\node[box, fill=red!8, right=0.5cm of head] (out) {Logits\\{\footnotesize \code{[B,S,V]}}};
\node[box, fill=green!8, above=0.7cm of inj] (act) {Activation\\{\footnotesize \code{[B,3584]}}};

\draw[arr] (tok) -- (emb);
\draw[arr] (emb) -- (inj);
\draw[arr] (inj) -- (tr);
\draw[arr] (tr) -- (head);
\draw[arr] (head) -- (out);
\draw[arr] (act) -- (inj);
\end{tikzpicture}
\end{center}
\end{frame}


% ============================================================
% Section 2: Training Method
% ============================================================

\begin{frame}{AV / Training：SFT + RL}
\small

AV 训练分两阶段：先用 API 生成的 explanation 做 warm-start，再用 RL 以重建保真度为 reward 继续优化。

\begin{columns}[T,onlytextwidth]
\begin{column}{0.48\linewidth}
\begin{block}{Stage 1a: AV-SFT（监督微调）}
\begin{itemize}
  \item 用 GPT-4 / Claude 等 API 为 activation 生成 reference explanation。
  \item 训练 AV 模彷 API 的解释风格和格式。
  \item 数据格式：prompt 含 \code{<INJECT>} token，response 为 \code{<explanation>...\\</explanation>}。
  \item Warm-start 让 AV 学会基本的 explanation 格式，为后续 RL 提供合理的初始策略。
\end{itemize}
\end{block}
\end{column}
\begin{column}{0.48\linewidth}
\begin{block}{Stage 1b: AV-RL（强化学习）}
\begin{itemize}
  \item 用 \textbf{GRPO} 继续训练。
  \item Reward = \textbf{AR cosine similarity} + \textbf{brevity bonus}。
  \item 鼓励 AR 能精确重建的紧凑 explanation。
\end{itemize}
\end{block}
\end{column}
\end{columns}

\vspace{0.5em}

Reward 信号来自\textbf{重建保真度}（AR cosine similarity）。但实践中，explanation 随 RL 训练自然地变得更 informative——因为更精确的描述有助于 AR 重建。

\end{frame}